% AUTO-GENERATED by the update_record tool. DO NOT EDIT.
% verify_paper regenerates this file from record.json and the run outputs
% and fails on any difference, so manual edits always surface.
\noindent\textbf{Hypothesis 1.} SciGym 公式の ReAct エージェントを弱い API モデル（gpt-4.1-mini / gpt-4.1-nano）で動かしても、論文の「上位モデルほど STE が低い」という順序が保たれ、gpt-4.1-mini の STE は論文の値から大きく悪化しない。~\cite[s1.p4, s1.p5]{duan-2025-measuring}
\begin{enumerate}
\item[\textbf{Claim 1}] SciGym-small の同じ 3 件で、gpt-4.1-nano の ReAct は gpt-4.1-mini の ReAct より STE が低くない。~\cite[s1.p1, s1.p2, s1.p3, s1.p4, s1.p7]{duan-2025-measuring}
  \emph{Rationale:} 論文の上位モデルほどよいという順序（s1.p4）を、さらに下位の nano と mini の対で確かめる。h1 の前半を担う。
  \emph{Criterion:} \texttt{\detokenize{react-gpt-4.1-nano}}.\texttt{\detokenize{trajectory_smape}} $-$ \texttt{\detokenize{react-gpt-4.1-mini}}.\texttt{\detokenize{trajectory_smape}} $\geq 0$.
  \emph{Prediction:} $[0, 0.2]$ (論文 Table 1 の上位版と mini 版の STE の差（GPT 0.4611 vs 0.6007 で 0.14、Gemini 0.3212 vs 0.4181 で 0.10、s1.p4）。3 件なので幅を広く取る).
  \emph{Observed:} 0.0992214.
  \emph{Verdict:} supported.
\item[\textbf{Claim 2}] SciGym-small の 3 件で、gpt-4.1-mini の ReAct の STE は論文 Table 1 の GPT-4.1-mini の値（0.6007）を 0.10 より大きくは上回らない。~\cite[s1.p1, s1.p2, s1.p3, s1.p5, s1.p6, s1.p7]{duan-2025-measuring}
  \emph{Rationale:} 同じ Controller・同じ設定で論文の数値に近い水準が出ることは、アダプタ（Vercel 経由の OpenAI 互換 LLM）が公式の手順を保っている証拠になり、h1 の後半を担う。c1 の参照 run でもある。
  \emph{Criterion:} \texttt{\detokenize{react-gpt-4.1-mini}}.\texttt{\detokenize{trajectory_smape}} $-$ 0.6007 $\leq 0.1$.
  \emph{Prediction:} $[-0.15, 0.1]$ (論文 Table 1 の GPT-4.1-mini の平均 STE 0.6007（s1.p5、137 件の平均）。3 件の平均は件の選び方で大きく振れるので幅を広く取る).
  \emph{Observed:} -0.22952.
  \emph{Verdict:} supported.
\end{enumerate}
\noindent\emph{Assumed, so that the claims together imply Hypothesis 1:}
\begin{itemize}
\item STE（trajectory\_smape、SciGym 公式 Evaluator、airas-eval scigym\_small）が機構回復の良さを表す（c1, c2）
\item 上流コード h4duan/SciGym @88a7b93 は論文の Controller を実装している（c1, c2）
\item knob の一覧（design の arguments）は、コードから列挙できた範囲を超えて完全である（c1, c2）
\item Vercel AI Gateway が返すモデル（openai/gpt-4.1-mini、openai/gpt-4.1-nano）はその名前のモデルであり、gpt-4.1-mini は論文の GPT-4.1-Mini-2025-04-14 と同じ重みである（c1, c2）
\item LLM のサンプリングは非決定的であり、1 回の run の差は再現時にずれうる（c1, c2）
\item Seyval（RIKYU の BYO Slurm）の実行記録と git 履歴は正しい（c1, c2）
\item 3 件は 137 件の代表ではない。pilot の目的は record の運用の確認であり、3 件での差が母集団の差を反映するとみなす（c1, c2）
\end{itemize}
